\section{Variabili Aleatorie}

Quando lo spazio campionario $\Omega$ non è un insieme numerico, è necessario mettere gli eventi in corrispondenza con un insieme numerico discreto. Come facciamo? Tramite una corrispondenza (funzione) definita come:

\[
X: \omega \in \Omega \to X(\omega) \in \mathcal{X} \subseteq \mathbb{R}
\]

Questa funzione $X$ prende il nome di \textbf{Variabile Aleatoria}.
\\Il codominio di $X$ (che nel caso discreto deve essere un insieme numerabile) prende il nome di \textbf{alfabeto} di $X$.

Osserviamo che:
\begin{itemize}
    \item Gli eventi elementari divengono del tipo $X(\omega) = x$, con $x \in \mathcal{X}$.
    \item Di conseguenza, anche le frequenze e le probabilità saranno espresse in funzione di $x$: $f_n(x)$ e $\mathbb{P}(X=x)$. La scrittura $\mathbb{P}(X=x)$ é una scorciatoia per: $P({\omega \in \Omega \mid X(\omega) = x})$, stiamo dicendo: Consideriamo tutti i casi in cui il risultato dell'esperimento (X) coincide con il numero (x) e calcoliamone la probabilitá. 
\end{itemize}

\paragraph{Etimologia e Terminologia}
La variabile aleatoria è anche detta variabile \textit{casuale} o \textit{stocastica}.
\begin{itemize}
    \item \textbf{Stocastico:} Deriva dal verbo greco \textit{stochazesthai} che significa "fare congetture" o "tirare al bersaglio". In origine il termine indicava qualcosa che si basava su ipotesi e congetture.
    \item \textbf{Aleatorio:} Deriva dal latino \textit{alea} (gioco di dadi) ed esprime il concetto di rischio calcolato.
\end{itemize}

\textbf{Definizione:} Ricordiamo che viene definito \textit{aleatorio} (o \textit{stocastico}) un fenomeno non deterministico.
\\Quindi una variabile casuale è una variabile che può assumere valori diversi in dipendenza da qualche fenomeno aleatorio.

\subsection{Funzione di Massa di Probabilità (PMF)}

In alcune situazioni potremmo essere interessati non alla probabilità che la variabile casuale $X$ assuma uno specifico valore, bensì alla probabilità che essa assuma un valore minore o uguale ad un dato valore $x$.

Sia $X$ una variabile aleatoria discreta con supporto (alfabeto) $\mathcal{X}$.
\\La funzione $p_X(x)$ definita come:

\begin{equation}
\boxed{p_X(x) = \mathbb{P}(X=x), \quad \forall x \in \mathcal{X}}
\end{equation}

è chiamata \textbf{Funzione di Massa di Probabilità} (pmf).
\\\textit{(Nota: talvolta indicata impropriamente come densità o distribuzione).}

Affinché questa funzione sia valida, deve soddisfare due proprietà assiomatiche che derivano dall'approccio frequentista (dove la probabilità è approssimata dal limite delle frequenze relative $\approx \lim_{n \to \infty} \frac{n_x}{n}$):

\begin{enumerate}
    \item \textbf{Non negatività:}
    \[ p_X(x) \ge 0, \quad \forall x \in \mathcal{X} \]
    
    \item \textbf{Normalizzazione:}
    \[ \sum_{x \in \mathcal{X}} p_X(x) = 1 \]
\end{enumerate}

\paragraph{Esempio di verifica}
Sia $\mathcal{X} = \{1, 2, 3, 4\}$. La sequenza $\left(\frac{1}{2}, \frac{1}{4}, \frac{1}{8}, \frac{1}{8}\right)$ definisce una PMF valida?

\begin{itemize}
    \item \textbf{Verifica 1:} Tutti i valori sono $\ge 0$? \textbf{Sì}.
    \item \textbf{Verifica 2:} La somma è unitaria?
    \[
    \sum_{i=1}^4 p_X(x_i) = \frac{1}{2} + \frac{1}{4} + \frac{1}{8} + \frac{1}{8} = \frac{4+2+1+1}{8} = \frac{8}{8} = 1
    \]
    \item \textbf{Conclusione:} Sì, è una PMF valida.
\end{itemize}

\subsection{Caratterizzazione e Media Campionaria}

Il modo più preciso per conoscere una variabile aleatoria è avere la sua \textbf{pmf} (Probability Mass Function), ovvero la lista completa di tutti i possibili valori che può assumere e le relative probabilità. Se hai questo, sai tutto della variabile.

Spesso non possiamo (o non vogliamo) elencare tutte le probabilità. Ci serve un "riassunto". Uno dei riassunti più utili è la \textbf{Media}.

\paragraph{L'Esperimento (prove)}
L'approccio empirico (statistico) si basa sui seguenti passi:
\begin{itemize}
    \item Immagina di avere la tua variabile $X$ (es. lanciare un dado o misurare un pezzo meccanico).
    \item Ripeti l'esperimento $n$ \textbf{volte}.
    \item Ottieni una collezione di $n$ risultati:
    \[ [X(\omega_1), X(\omega_2), \dots, X(\omega_n)] \]
\end{itemize}

\paragraph{La Media Campionaria ($\overline{X}_n$)}
Per capire come si comporta la variabile "in media", calcoli la media aritmetica dei risultati che hai ottenuto. La formula è:

\begin{equation}
\boxed{\overline{X}_n = \frac{1}{n} \sum_{i=1}^n X(\omega_i)}
\end{equation}

Ovvero: \textbf{Somma di tutti i risultati ottenuti diviso il numero di prove fatte.}\

\subsection{Media Statistica (Valore Atteso)}

Riconsideriamo la definizione di media campionaria vista in precedenza:
\[
\overline{X}_n = \frac{1}{n} \sum_{i=1}^n X(\omega_i)
\]

Supponiamo che la variabile aleatoria $X$ assuma valori in un insieme discreto finito (il supporto) $\mathcal{X} = \{x_1, \dots, x_M\}$.
\\Al crescere del numero di esperimenti $n$, accadrà che il valore specifico $x_1$ verrà osservato $n_1$ volte, il valore $x_2$ verrà osservato $n_2$ volte, e così via (fino a $x_M$ che apparirà $n_M$ volte).

Possiamo quindi riscrivere la somma non più come lista di tutti i risultati, ma raggruppando i valori uguali e moltiplicandoli per la loro frequenza di apparizione ($n_i$).
\\Quando il numero di prove $n$ tende all'infinito ($n \to \infty$), le frequenze relative $f_n(x_i) = \frac{n_i}{n}$ convergono alle probabilità teoriche $p_X(x_i)$.

Questo ragionamento ci porta alla definizione teorica:

\begin{equation}
\begin{aligned}
\overline{X}_n &= \frac{1}{n} \sum_{i=1}^M n_i x_i \\
&= \sum_{i=1}^M x_i \underbrace{\left( \frac{n_i}{n} \right)}_{f_n(x_i)} \\
&\xrightarrow{n \to \infty} \sum_{i=1}^M x_i \mathbb{P}(X = x_i) \\
&= \sum_{i=1}^M x_i p_X(x_i) \stackrel{\text{def}}{=} \mathbb{E}[X]
\end{aligned}
\end{equation}

\paragraph{Definizione}
La quantità limite ottenuta è definita come \textbf{Media Statistica} o \textbf{Valore Atteso} (Expected Value) della variabile aleatoria $X$ ed è indicata con la notazione $\mathbb{E}[X]$:

\begin{equation}
\boxed{\mathbb{E}[X] = \sum_{x \in \mathcal{X}} x \cdot p_X(x)}
\end{equation}
in pratica é come una media pesata dove (x) é il peso.

\textit{Nota: Nel caso in cui il supporto $\mathcal{X}$ sia infinito numerabile ($M=\infty$), la somma diventa una serie e il concetto si estende come passaggio al limite.}

\subsection{Varianza e Deviazione Standard}
\label{sec-varianza}

Il Valore Atteso (media) ci fornisce il "baricentro" della distribuzione, ma non ci dice nulla su quanto i dati siano sparsi attorno ad esso.
\\Per descrivere la dispersione dei valori (ovvero quanto è "incerta" o "rischiosa" la variabile aleatoria), introduciamo il concetto di \textbf{Varianza}.

\paragraph{}
Misuriamo la "distanza" tra ogni possibile valore ($X$) e la media $\mu = \mathbb{E}[X]$. Dove per distanza intendiamo quanto é lontano $x$ da $\mu$ 
\begin{itemize}
    \item La differenza $X - \mu$ è detta \textbf{scarto}.
    
    \item Se facessimo la media degli scarti ($\mathbb{E}[X - \mu]$), otterremmo sempre 0 (gli scarti positivi cancellano quelli negativi). 
    
    \begin{center}
    \fbox{
    \begin{minipage}{0.9\linewidth}
    \small \textbf{Perché fa sempre 0?} \\
    Partendo dalla definizione di media:
    $$ \mu = \frac{x_1 + x_2 + \dots + x_n}{n} $$
    Moltiplicando entrambi i lati per $n$, la somma di tutti i valori originali equivale alla media presa $n$ volte:
    $$ x_1 + x_2 + \dots + x_n = n \cdot \mu $$
    Se ora sommiamo tutti gli scarti:
    $$ (x_1 - \mu) + (x_2 - \mu) + \dots + (x_n - \mu) $$
    Possiamo riordinare e raggruppare i termini:
    $$ (x_1 + x_2 + \dots + x_n) - (\mu + \mu + \dots + \mu) $$
    Sapendo che la prima parentesi vale $n \cdot \mu$, e la seconda contiene $\mu$ sommata $n$ volte:
    $$ n \cdot \mu - n \cdot \mu = 0 $$
    \end{minipage}
    }
    \end{center}
    
    \item Per evitare ciò, eleviamo gli scarti al quadrato.
\end{itemize}
\paragraph{Definizione}
La Varianza di una variabile aleatoria $X$, indicata con $\text{Var}(X)$ o $\sigma_X^2$, è definita come la \textbf{media degli scarti al quadrato}:

\begin{equation}
\boxed{\text{Var}(X) = \mathbb{E}\left[ (X - \mathbb{E}[X])^2 \right]}
\end{equation}

Essendo una somma di quadrati pesata per delle probabilità (che sono positive), la varianza è sempre una quantità \textbf{non negativa} ($\text{Var}(X) \ge 0$).

\subsubsection{Come si calcola?}

Per calcolare la varianza negli esercizi, è spesso scomodo usare la definizione sopra (bisognerebbe calcolare ogni scarto, farne il quadrato e moltiplicare per la probabilità).
\\Sviluppando il quadrato del binomio all'interno del valore atteso, si ottiene una formula molto più pratica, nota come \textit{formula di König-Huygens}:

$$ \text{Var}(X) = \mathbb{E}[(X - \mu)^2] $$
$$ = \mathbb{E}[X^2 - 2X\mu + \mu^2] $$

Per la proprietà di linearità del valore atteso:
$$ = \mathbb{E}[X^2] - \mathbb{E}[2X\mu] + \mathbb{E}[\mu^2] $$

Poiché $2$ e $\mu$ sono costanti, possono essere portate fuori dall'operatore $\mathbb{E}$ (e il valore atteso della costante $\mu^2$ è $\mu^2$ stesso):
$$ = \mathbb{E}[X^2] - 2\mu \mathbb{E}[X] + \mu^2 $$
$$ = \mathbb{E}[X^2] - 2\mu \cdot \mu + \mu^2 $$
$$ = \mathbb{E}[X^2] - 2\mu^2 + \mu^2 $$
$$ = \mathbb{E}[X^2] - \mu^2 $$

\begin{equation}
\boxed{\text{Var}(X) = \mathbb{E}[X^2] - (\mathbb{E}[X])^2}
\end{equation}

In parole povere: \textbf{"La media dei quadrati meno il quadrato della media"}.
Dove il termine $\mathbb{E}[X^2]$ (detto \textit{momento secondo o valore quadratico medio (Mean Square)}) si calcola come:
\[ \mathbb{E}[X^2] = \sum_{x \in \mathcal{X}} x^2 \cdot p_X(x) \]

\paragraph{Esempio di Calcolo}
Consideriamo la variabile aleatoria $X$ definita come il lancio di una moneta dove Testa=1 e Croce=0, con probabilità $1/2$.
\begin{enumerate}
    \item Calcoliamo la media $\mathbb{E}[X] = 0 \cdot 0.5 + 1 \cdot 0.5 = 0.5$.
    \item Calcoliamo la media dei quadrati $\mathbb{E}[X^2] = 0^2 \cdot 0.5 + 1^2 \cdot 0.5 = 0.5$.
    \item Calcoliamo la varianza:
    \[ \text{Var}(X) = 0.5 - (0.5)^2 = 0.5 - 0.25 = 0.25 \]
\end{enumerate}
Ma che significa 0.25 Testa al quadrato? Facciamone la radice quadrata: $\sqrt{0.25}=0.5$ Testa. In questo caso specifico, non importa cosa esce (testa o croce), ti sbaglierai sempre esattamente di 0.5 rispetto alla media. 

\subsubsection{Deviazione Standard (Scarto Quadratico Medio)}

La varianza ha un difetto "fisico": la sua unità di misura è al quadrato rispetto alla variabile originale (es. se $X$ è in metri, la varianza è in $m^2$).
\\Per tornare all'unità di misura originale e avere un indice di dispersione più intuitivo, si definisce la \textbf{Deviazione Standard} ($\sigma_X$):

\begin{equation}
\sigma_X = \sqrt{\text{Var}(X)}
\end{equation}

Mentre la varianza è utile per i calcoli algebrici (grazie alle proprietà additive), la deviazione standard è fondamentale per l'interpretazione dei dati.

\section{Relazione tra Misura, Variabile Aleatoria e Distribuzione}

Per fare chiarezza su quanto detto, analizziamo le differenze e le relazioni tra i concetti chiave.

\paragraph{Definizioni Fondamentali}
\begin{itemize}
    \item \textbf{La Misura di Probabilità ($\mathbb{P}$):} È una funzione che prende in ingresso un \textit{evento} (un insieme di risultati) e restituisce un numero reale compreso tra 0 e 1. Essa quantifica quanto è probabile che accada quel certo evento.
    
    \item \textbf{La Variabile Aleatoria ($X$):} È una funzione che prende un singolo risultato astratto ($\omega \in \Omega$) e lo trasforma in un numero reale.
    \\Serve a tradurre un risultato "fisico" o astratto in un valore quantitativo.
    \\\textit{Esempio:} "Esce la faccia del Re di cuori" $\xrightarrow{X}$ "vale 10 punti".
\end{itemize}

\subsection{La Distribuzione (Misura Indotta)}

La misura di probabilità indotta sullo spazio misurabile di arrivo $(E, \mathcal{F})$ da una variabile aleatoria $X$, a partire dalla misura originale $\mathbb{P}$ su $(\Omega, \mathcal{E})$, è detta \textbf{distribuzione}, \textbf{legge} o \textbf{probabilità di X}, ed è indicata con $P_X$.

In pratica, è la misura associata alla variabile aleatoria. Mentre la misura originale $\mathbb{P}$ agisce sugli eventi dello spazio campionario, la distribuzione $P_X$ restituisce la probabilità dei valori numerici in output.

\textbf{Concetto chiave:}
\begin{quote}
    È la misura $\mathbb{P}$ che è stata \textbf{"trasportata"} sui numeri reali tramite la variabile aleatoria $X$.
\end{quote}


Ovviamente, la forma e il comportamento della distribuzione $P_X$ dipendono interamente dalla definizione della variabile aleatoria $X$.

\subsection{Tipologie di Distribuzioni}

Le distribuzioni si dividono in due grandi famiglie in base alla natura del loro codominio:

\begin{enumerate}
    \item \textbf{Discrete:} Se la funzione $X$ mappa i risultati su un insieme di numeri che si possono \textit{contare} (insieme numerabile).
    \item \textbf{Continue:} Se il codominio è un intervallo continuo di numeri reali.
\end{enumerate}

A differenza delle variabili continue (dove la probabilità in un singolo punto esatto è matematicamente zero), nelle distribuzioni discrete la probabilità è concentrata in \textbf{"grumi"} o \textbf{"masse"} su specifici punti.

La funzione di massa di probabilità $p_X(x)$ risponde alla domanda:
\[
\text{"\textbf{Esattamente quanta probabilità è accumulata su questo specifico numero?}"}
\]


\subsection{Esempio: Il Conteggio Bernoulliano}

Per applicare concretamente i concetti appena esposti, consideriamo un esperimento fondamentale nella teoria della probabilità: il lancio ripetuto di una moneta. Questo scenario è il prototipo del cosiddetto \textbf{Processo di Bernoulli}.

\paragraph{Descrizione dell'Esperimento}
Immaginiamo di lanciare una moneta per $N$ volte consecutive. Possiamo modellare questo fenomeno attraverso tre ipotesi chiave:
\begin{itemize}
    \item \textbf{Esiti binari:} Ogni singolo lancio ammette solo due risultati: "Testa" ($T$), che chiameremo successo, e "Croce" ($C$), che chiameremo insuccesso.
    \item \textbf{Probabilità costanti:} La moneta ha una probabilità $p$ (con $0 < p < 1$) di dare "Testa". Di conseguenza, la probabilità di ottenere "Croce" sarà $q = 1 - p$.
    \item \textbf{Indipendenza:} Assumiamo che i lanci siano statisticamente indipendenti; l'esito di un lancio non influenza in alcun modo quelli precedenti o successivi.
\end{itemize}

\paragraph{Ora Passiamo dallo Spazio Campione alla Variabile Aleatoria}
Lo spazio campione $\Omega$ è costituito dall'insieme di tutte le possibili sequenze di risultati (le $N$-ple). Un evento elementare generico $\omega \in \Omega$ si presenta come una stringa ordinata del tipo:
\[ \omega = (C, C, T, T, C, \dots, C, T, T) \]

Il nostro obiettivo non è analizzare l'ordine esatto, ma il numero complessivo di successi. Definiamo quindi la variabile aleatoria $X_N$ che "conta" quante volte esce "Testa" negli $N$ lanci:
\[ X_N: \omega \in \Omega \longrightarrow X_N(\omega) \in \{0, 1, \dots, N\} \]
Diremo che $X_N(\omega) = k$ se la sequenza $\omega$ contiene esattamente $k$ successi.

\subsubsection{Costruzione della Probabilità}

Per determinare la probabilità che si verifichino esattamente $k$ successi ($P(X_N=k)$), procediamo per gradi sfruttando l'indipendenza degli eventi.

\paragraph{1. Analisi di una singola configurazione}
Consideriamo una specifica sequenza $\omega^*$ che soddisfa la condizione di avere $k$ teste e $N-k$ croci. Per fissare le idee, immaginiamo la sequenza in cui appaiono prima tutte le teste e poi tutte le croci. Grazie all'indipendenza stocastica, la probabilità congiunta è semplicemente il prodotto delle singole probabilità:

\[
\mathbb{P}(\omega^*) = \mathbb{P}(\underbrace{T \cap \dots \cap T}_{k} \cap \underbrace{C \cap \dots \cap C}_{N-k}) = p^k q^{N-k}
\]
Questa rappresenta la probabilità di una singola, specifica sequenza in cui i risultati si presentano in un ordine ben preciso.
È fondamentale notare che qualsiasi altra sequenza con lo stesso numero di teste e croci (anche se in ordine diverso) avrà la \textbf{stessa identica probabilità}, poiché il prodotto è commutativo.

\paragraph{2. Il ruolo della Combinatoria}
Poiché gli eventi elementari sono mutuamente esclusivi, la probabilità totale dell'evento "ottenere $k$ teste" è la somma delle probabilità di tutte le sequenze favorevoli.
\\Dobbiamo quindi chiederci: \textit{quante sono le sequenze distinte di lunghezza $N$ che contengono esattamente $k$ teste?}
\\La risposta è fornita dal \textbf{coefficiente binomiale}:
\[ C_{N,k} = \binom{N}{k} \]
L'insieme di queste sequenze costituisce l'evento composto $\Omega_{N,k}$.

\subsection{La Distribuzione Binomiale}

Unendo il conteggio combinatorio con la probabilità della singola sequenza, otteniamo la \textbf{Funzione di Massa di Probabilità} (PMF) Binomiale.
\\La probabilità che la variabile aleatoria $X_N$ assuma valore $k$ è data da:

\begin{equation}
\boxed{p_{X_N}(k) = \mathbb{P}(X_N = k) = \binom{N}{k} p^k q^{N-k}}
\end{equation}

Quando una variabile segue questa legge, si usa la notazione sintetica $X_N \sim \mathcal{B}(N, p)$.

\paragraph{Verifica di Normalizzazione}
Una proprietà essenziale di ogni PMF è che la somma delle probabilità su tutto l'alfabeto sia pari a 1. Nel nostro caso, sfruttando il celebre \textbf{Teorema del Binomio di Newton}, osserviamo che:

\[
\sum_{k=0}^{N} p_{X_N}(k) = \sum_{k=0}^{N} \binom{N}{k} p^k q^{N-k} = (p + q)^N
\]

Poiché $p$ e $q$ sono complementari ($p+q=1$), otteniamo $(1)^N = 1$, confermando che la distribuzione è ben definita.

\paragraph{Valore Atteso}
Infine, calcoliamo la media statistica della distribuzione. Applicando la definizione di valore atteso, si dimostra che:

\begin{equation}
\mathbb{E}[X_N] = \sum_{k=0}^{N} k \cdot \binom{N}{k} p^k q^{N-k} = Np
\end{equation}

Il risultato $\mathbb{E}[X_N] = Np$ è coerente con l'intuizione: se ripetiamo un esperimento $N$ volte con probabilità di successo $p$, ci aspettiamo in media $N \cdot p$ successi.

\subsection{La Distribuzione Uniforme Discreta}

Se la variabile di Bernoulli modellava una situazione con esattamente 2 esiti, la \textbf{Variabile Uniforme} modella situazioni in cui tutti gli esiti possibili hanno la \textbf{stessa probabilità} di verificarsi (sono equiprobabili).
\\Questa distribuzione concerne tutti gli esempi classici visti fino ad ora, come il lancio di un dado onesto, l'estrazione di una carta da un mazzo completo, ecc.

\paragraph{Definizioni}
Una variabile aleatoria $X$ si dice \textit{uniformemente distribuita} su un insieme (alfabeto) $\mathcal{X}$ (e si indica con $X \sim \mathcal{U}(\mathcal{X})$) se:
\begin{itemize}
    \item L'insieme dei valori possibili ha cardinalità finita $M$ (ossia $|\mathcal{X}| = M$).
    \item Ogni singolo valore $x \in \mathcal{X}$ ha probabilità $1/M$.
\end{itemize}

La sua funzione di massa di probabilità è quindi definita come:

\begin{equation}
\boxed{p_X(x) = \mathbb{P}(X=x) = \frac{1}{M}, \quad \forall x \in \mathcal{X}}
\end{equation}

Ovviamente $p_X(x)$ soddisfa le condizioni necessarie per essere una pmf (valori positivi e somma unitaria).

\paragraph{Esempio Classico}
Consideriamo il lancio di un dado a 6 facce onesto.
\begin{itemize}
    \item \textbf{Alfabeto:} $\mathcal{X} = \{1, 2, 3, 4, 5, 6\}$.
    \item \textbf{Cardinalità:} $M = 6$.
    \item \textbf{Probabilità puntuale:} La probabilità che esca un numero specifico (es. il 4) è:
    \[ \mathbb{P}(X=4) = \frac{1}{6} \]
\end{itemize}

\subsubsection{Valore Atteso (Media)}

Il calcolo della media è immediato. Applicando la definizione generale:

\begin{equation}
\mathbb{E}[X] = \sum_{x \in \mathcal{X}} x \cdot p_X(x) = \sum_{x \in \mathcal{X}} x \cdot \frac{1}{M} = \frac{1}{M} \sum_{x \in \mathcal{X}} x
\end{equation}

Poiché la probabilità $1/M$ è costante, la portiamo fuori dalla sommatoria. È ovvio notare che, in questo caso specifico, la media della variabile aleatoria \textbf{coincide con la media aritmetica} di tutti i possibili valori dell'alfabeto.

\subsubsection{Un Caso Notevole: Interi consecutivi da 0}

Un caso molto interessante, frequente in informatica, si ha quando l'alfabeto è costituito dai numeri interi da $0$ a $M-1$.
\[ \mathcal{X} = \{0, 1, \dots, M-1\} \]

Che succede al valore atteso in questo caso? Per calcolarlo sfruttiamo la \textbf{formula della somma dei primi interi} (sommatoria di Gauss).

\begin{enumerate}
    \item Dobbiamo sommare i numeri da $0$ a $M-1$.
    \item Ricordiamo che la somma da $1$ a $n$ è data da $\frac{n(n+1)}{2}$.
    \item Nel nostro caso il numero più alto è $n = M-1$.
    \item Sostituendo nella formula della somma:
    \[ \text{Somma} = \sum_{i=0}^{M-1} i = \frac{(M-1)((M-1)+1)}{2} = \frac{(M-1)M}{2} \]
    \item Ora calcoliamo la media dividendo per $M$ (come richiesto dalla formula della media uniforme):
    
    \begin{equation}
    \begin{aligned}
    \mathbb{E}[X] &= \frac{1}{M} \cdot \text{Somma} \\
    &= \frac{1}{M} \cdot \frac{M(M-1)}{2}
    \end{aligned}
    \end{equation}
    
    \item Semplificando $M$, otteniamo la formula finale per questo caso specifico:
    
    \begin{equation}
    \boxed{\mathbb{E}[X] = \frac{M-1}{2}}
    \end{equation}
\end{enumerate}

\paragraph{Esempio numerico}
Ipotizziamo di avere una cifra decimale generata a caso (da 0 a 9). In questo caso $M=10$.
\begin{itemize}
    \item I valori sono $\{0, 1, \dots, 9\}$.
    \item Applicando la formula appena ricavata:
    \[ \mathbb{E}[X] = \frac{10-1}{2} = \frac{9}{2} = 4.5 \]
    \item Il valore atteso è $4.5$.
\end{itemize}

\subsection{La Distribuzione di Poisson}

Prima di definire formalmente questa distribuzione, consideriamo uno scenario intuitivo per comprenderne l'utilità.

\paragraph{Esempio: Il Call Center}
Immaginiamo di lavorare in un call center. Analizzando i dati storici, sappiamo che, \textbf{in media}, riceviamo \textbf{5 chiamate all'ora}.
\\Tuttavia, il fatto che la media sia 5 non garantisce che nella prossima ora riceveremo \textit{esattamente} 5 chiamate. Essendo un fenomeno aleatorio:
\begin{itemize}
    \item In un'ora potremmo riceverne 3 (bassa attività).
    \item In un'altra potremmo riceverne 8 (picco di attività).
    \item In un'altra ancora potremmo non riceverne nessuna (0).
\end{itemize}

La distribuzione di Poisson nasce proprio per rispondere a domande specifiche su questi conteggi, del tipo:
\begin{quote}
    \textit{"Sapendo che la media è 5 chiamate/ora, qual è la probabilità che nella prossima ora ne arrivino esattamente 7?"}
\end{quote}

È la distribuzione di riferimento per modellare il conteggio di \textbf{eventi rari} che si verificano in un intervallo continuo di tempo, spazio o volume.

\paragraph{Il Parametro $\lambda$ e le Condizioni}
Tutta la distribuzione ruota attorno a un solo parametro fondamentale, $\lambda$ (Lambda), che rappresenta l'\textbf{intensità media} del processo:

\[ \lambda = \text{numero medio di eventi attesi nell'intervallo} \]

Affinché si possa applicare il modello di Poisson, devono valere tre ipotesi:
\begin{enumerate}
    \item \textbf{Indipendenza:} L'accadimento di un evento non influenza la probabilità che ne accada un altro (es. ricevere una chiamata ora non cambia la probabilità di riceverne una tra un minuto).
    \item \textbf{Stazionarietà:} Il tasso medio $\lambda$ è costante nel tempo considerato.
    \item \textbf{Non contemporaneità:} La probabilità che due eventi accadano nello stesso identico istante è trascurabile.
\end{enumerate}

\subsubsection{Derivazione come limite della Binomiale}

La formula di Poisson non è arbitraria, ma si ottiene matematicamente come caso limite della distribuzione Binomiale.

Immaginiamo di suddividere il nostro intervallo di tempo (l'ora del call center) in $N$ sotto-intervalli piccolissimi (es. millisecondi).
\begin{itemize}
    \item Sia $N \to \infty$ (numero di sotto-intervalli tendente all'infinito).
    \item Sia $p \to 0$ (probabilità di ricevere una chiamata nel singolo millisecondo è quasi nulla).
    \item Manteniamo però costante il prodotto $N \cdot p = \lambda$ (la media totale delle chiamate nell'ora resta 5).
\end{itemize}

Partiamo dalla PMF Binomiale e sostituiamo $p$ con $\frac{\lambda}{N}$:
N.B: ${p} = \frac{\lambda}{N}$
\[
\mathbb{P}(X=k) = \lim_{N \to \infty} \binom{N}{k} \left( \frac{\lambda}{N} \right)^k \left( 1 - \frac{\lambda}{N} \right)^{N-k}
\]

Espandiamo i termini per calcolare il limite:

\begin{equation}
\begin{aligned}
\mathbb{P}(X=k) &= \lim_{N \to \infty} \frac{N!}{k!(N-k)!} \cdot \frac{\lambda^k}{N^k} \cdot \left( 1 - \frac{\lambda}{N} \right)^N \cdot \left( 1 - \frac{\lambda}{N} \right)^{-k} \\
&= \frac{\lambda^k}{k!} \lim_{N \to \infty} \underbrace{\frac{N(N-1)\dots(N-k+1)}{N^k}}_{\to 1} \cdot \underbrace{\left( 1 - \frac{\lambda}{N} \right)^N}_{\to e^{-\lambda}} \cdot \underbrace{\left( 1 - \frac{\lambda}{N} \right)^{-k}}_{\to 1}
\end{aligned}
\end{equation}

Dove abbiamo sfruttato il limite notevole $\lim_{x \to \infty} (1 - \frac{\lambda}{x})^x = e^{-\lambda}$.

\paragraph{Funzione di Massa di Probabilità}
Il risultato del limite ci fornisce la formula della PMF per una variabile Poissoniana $X \sim \text{Pois}(\lambda)$:

\begin{equation}
\boxed{p_X(k) = \mathbb{P}(X=k) = \frac{\lambda^k \cdot e^{-\lambda}}{k!}, \quad k \in \{0, 1, 2, \dots\}}
\end{equation}

Dove $e$ è il numero di Nepero ($e \approx 2.718$).

\paragraph{Proprietà: Media e Varianza}
Una caratteristica unica ed elegante della distribuzione di Poisson è che il valore atteso e la varianza coincidono e sono pari al parametro stesso.

\begin{equation}
\mathbb{E}[X] = \lambda \quad \text{e} \quad \text{Var}(X) = \lambda
\end{equation}

Questo implica che in un processo di Poisson, se aumenta il numero medio di eventi attesi (es. un call center più grande), aumenta proporzionalmente anche l'incertezza (la variabilità) attorno a tale media.

\section{Pmf e medie condizionali}

Supponiamo che $X \sim \mathcal{B}(16, 1/3)$.

\textbf{"La variabile aleatoria $X$ segue una distribuzione Binomiale con parametri $n= 16$ e $p=1/3$"}

dove $16$ è il numero di volte che ripeti l'esperimento e $1/3$ è la probabilità di vincere nel singolo esperimento.

$X$ è il risultato finale: \textbf{conta quante volte hai vinto} su quei $16$ tentativi. Il minimo che $X$ può assumere è $0$ (hai perso tutte e $16$ le volte) e il massimo che $X$ può assumere è $16$ (hai vinto sempre).

Per definizione di distribuzione binomiale:

\begin{equation}
\boxed{p_{X_N}(k) = \mathbb{P}(X_N = k) = \binom{N}{k} p^k q^{N-k}}
\end{equation}

quindi:

\[
\binom{16}{k} \left(\frac{1}{3}\right)^k \left(\frac{2}{3}\right)^{16-k}
\]

Significa:
\begin{enumerate}
    \item $\left(\frac{1}{3}\right)^k$: Hai vinto $k$ volte.
    \item $\left(\frac{2}{3}\right)^{16-k}$: Hai perso tutte le altre volte.
    \item $\binom{16}{k}$: È il coefficiente binomiale. Ti dice in quanti modi diversi potevi ottenere quelle vittorie 
\end{enumerate}

Osserviamo che:
\[ \mathbb{E}[X] = 16/3 \quad \text{perché } \mathbb{E}[X] = N \cdot p \]
Se ripetiamo un esperimento $N$ volte con probabilità di successo $p$, ci aspettiamo in media $N \cdot p$ successi.

\paragraph{Scenario Condizionale}
Supponiamo ora di assumere che sia verificata la condizione $X > 4$: ci chiediamo come si distribuisca $X$ \textbf{sotto questa condizione}.

Dal punto di vista fisico significa considerare un insieme di $n$ esperimenti, scartare i valori di $X$ (il numero di successi) che siano inferiori a $5$ e valutare le probabilità dei residui dodici valori sul campione così ridotto.

Adesso devi ricalcolare le probabilità per il nuovo scenario "filtrato". Succedono due cose:

\paragraph{A. I valori scartati vanno a zero}

Per $k < 5$ (quindi $0, 1, 2, 3, 4$), la nuova probabilità è $0$.

\[ \text{per } k < 5 \to 0 \]

\paragraph{B. I valori rimasti devono "ingrassare" (Rinormalizzazione)}


Nella situazione originale, la somma delle probabilità dal $5$ al $16$ \textbf{non faceva 1} (perché mancava la "fetta" dello $0-4$). Diciamo, per esempio, che la probabilità di fare più di $4$ fosse il $60\%$ ($0.6$).

Se tu cancelli il resto, quel $60\%$ deve diventare il nuovo $100\%$ del tuo universo ridotto.
Per farlo, devi \textbf{dividere} ogni singola probabilità superstite per il totale della probabilità rimasta ($P(X > 4)$).

Quindi:
\begin{equation}
p_{X|X>4}(k) = \frac{p_X(k)}{P(X > 4)}
\end{equation}

\begin{itemize}
    \item \textbf{Numeratore ($p_X(k)$):} La probabilità originale di fare, ad esempio, $7$.
    \item \textbf{Denominatore ($P(X > 4)$):} La somma di tutte le probabilità "sopravvissute". Serve a scalare tutto in modo che la somma finale faccia di nuovo $1$.
\end{itemize}

La slide usa questa notazione formale:
\[
\frac{P(\{X = k\} \cap \{X > 4\})}{P(\{X > 4\})}
\]

Che é la stessa cosa che abbiamo appena visto 
\begin{itemize}
    \item \textbf{Caso $k = 3$:} L'intersezione tra "è uscito 3" e "il numero è maggiore di 4" è \textbf{impossibile} (insieme vuoto). Probabilità = $0$.
    \item \textbf{Caso $k = 7$:} L'intersezione tra "è uscito 7" e "il numero è maggiore di 4" è semplicemente \textbf{"è uscito 7"}. Quindi al numeratore metti la probabilità originale del $7$.
\end{itemize}

\paragraph{Proprietà della PMF condizionale}
Si noti che la pmf condizionale trovata è una vera pmf. Infatti:

\[
p_{X|X>4}(k) \ge 0 \quad \text{e} \quad \sum_{k \in \mathbb{Z}} p_{X|X>4}(k) = \sum_{k=5}^{16} \frac{p_X(k)}{\sum_{i=5}^{16} p_X(i)} = 1
\]

Siamo quindi ora in grado di definire la pmf di una variabile aleatoria qualsiasi $X$ condizionata a un qualsiasi evento $A \subseteq \Omega$ a probabilità non nulla nella forma:

\begin{equation}
\boxed{p_{X|A}(x) = \mathbb{P}(x|A) = \frac{\mathbb{P}(\{X = x\} \cap A)}{\mathbb{P}(A)} \quad x \in \mathcal{X}}
\end{equation}

Ovviamente, $p_{X|A}(x)$ al variare di $x$ in $\mathcal{X}$ con $A$ prefissato è una pmf.

\section{Regola della probabilità totale per le pmf}

Ricordiamo che, per un qualunque evento $C \subseteq \Omega$ e per una qualunque partizione $\{E_i\}_{i=1}^M$, si ha:

\[
\mathbb{P}(C) = \sum_{i=1}^M \mathbb{P}(C \mid E_i)\mathbb{P}(E_i)
\]

Ovvero: la probabilitá che l'evento C accada é uguale alla probabilità dell'evento C condizionata da ogni Evento indipendente Ei per la probabilità di ogni evento Ei.
\\\\
Pertanto, specializzando la precedente all'evento $C = \{X = x\}$, otteniamo una formula fondamentale per le funzioni di massa di probabilità:

\begin{equation}
\boxed{p_X(x) = \mathbb{P}(X = x) = \sum_{m=1}^M \mathbb{P}(\{X = x\} \mid E_m) \mathbb{P}(E_m) = \sum_{m=1}^M p_{X|E_m}(x)\mathbb{P}(E_m)}
\end{equation}

Il concetto chiave qui è la \textbf{Partizione} $(E_1, E_2, \dots, E_M)$.
\\Immagina una torta intera (l'universo $\Omega$). La tagli in tante fette ($E_m$). Le fette non si sovrappongono e coprono tutta la torta.
\\\\
Se vuoi sapere la probabilità totale che accada $X = x$ (es. che esca un 7), puoi procedere ricostruendo la probabilità "globale" usando i pezzi "locali":

\begin{enumerate}
    \item Calcola la probabilità che esca 7 nello \textbf{scenario 1} ($E_1$) e moltiplicala per quanto è probabile che accada lo scenario 1.
    \item Calcola la probabilità che esca 7 nello \textbf{scenario 2} ($E_2$) e moltiplicala per la probabilità dello scenario 2.
    \item ...fai così per tutti gli scenari...
    \item \textbf{Somma tutto.}
\end{enumerate}

\paragraph{Esempio applicato}
Tornando all'esempio precedente, vuoi calcolare $p_X(k)$ (probabilità di ottenere $k$ successi). Invece di calcolarla direttamente, spacca il mondo in due scenari:
\begin{itemize}
    \item \textbf{Scenario A:} $X > 4$ (hai vinto tante volte).
    \item \textbf{Scenario B:} $X \le 4$ (hai vinto poche volte).
\end{itemize}

La formula ti dice:
\[
P(\text{esce } k) = P(\text{esce } k \mid \text{scenario A}) \cdot P(A) + P(\text{esce } k \mid \text{scenario B}) \cdot P(B)
\]
\\
\paragraph{Formalmente:} Applicando la regola della probabilità totale, possiamo scrivere $p_X(k)$ in modi diversi a seconda di come decidiamo di partizionare lo spazio degli eventi.

Ecco due esempi di scomposizione validi per la stessa variabile:

\begin{align*}
p_X(k) &= p_{X \mid X>4}(k)\mathbb{P}(X > 4) + p_{X \mid X \le 4}(k)\mathbb{P}(X \le 4) \\[1em]
&= p_{X \mid 2 \le X \le 6}(k)\mathbb{P}(2 \le X \le 6) + p_{X \mid X \le 1}(k)\mathbb{P}(X \le 1) + p_{X \mid X \ge 7}(k)\mathbb{P}(X \ge 7)
\end{align*}

Queste uguaglianze sono valide \textbf{essendo} verificate le proprietà delle partizioni su $\Omega$:

\begin{itemize}
    \item \textbf{Caso a 2 scenari:}
    \[ \Omega = E_1 \cup E_2 = \{X > 4\} \cup \{X \le 4\} \]
    con $E_1 \cap E_2 = \emptyset$ (insiemi disgiunti).
    
    \item \textbf{Caso a 3 scenari:}
    \[ \Omega = E'_1 \cup E'_2 \cup E'_3 = \{2 \le X \le 6\} \cup \{X \le 1\} \cup \{X \ge 7\} \]
    con $E'_i \cap E'_j = \emptyset$ per ogni $i \neq j$ (tutti gli insiemi sono a due a due disgiunti).
\end{itemize}

\section{Medie e Medie Condizionali}

Qui si applica lo stesso identico ragionamento visto per le probabilità, ma al \textbf{Valore Atteso} (la Media). Questa è una delle formule più usate al mondo, spesso chiamata \textbf{Teorema della Media Totale}.

\paragraph{(Esempio della Scuola)}
Immagina di voler calcolare l'altezza media di \textbf{tutti} gli studenti di una scuola ($\mathbb{E}[X]$).
\\È difficile misurarli tutti insieme. È più facile fare così:

\begin{enumerate}
    \item Vai nella \textbf{Classe 1} ($E_1$), calcoli l'altezza media lì dentro ($\mathbb{E}[X \mid E_1]$) e la pesi per quanti alunni ha la classe.
    \item Vai nella \textbf{Classe 2} ($E_2$), calcoli la media lì ($\mathbb{E}[X \mid E_2]$) e pesi.
    \item Alla fine sommi i risultati pesati.
\end{enumerate}

\paragraph{La Formula}
Matematicamente, il teorema si esprime come:

\begin{equation}
\boxed{\mathbb{E}[X] = \sum_{m=1}^M \mathbb{P}(E_m) \cdot \mathbb{E}[X \mid E_m]}
\end{equation}

Dove:
\begin{itemize}
    \item $\mathbb{P}(E_m)$: Il "peso" dello scenario (es. quanto è grande la classe).
    \item $\mathbb{E}[X \mid E_m]$: La media calcolata \textit{solo} dentro quello scenario.
\end{itemize}

\subsection{Derivazione Matematica}

Partendo dalla definizione di valore atteso:

\[
\mathbb{E}[X] = \sum_{x \in \mathcal{X}} x p_X(x)
\]

Sostituiamo $p_X(x)$ con la regola della probabilità totale vista prima e scambiamo l'ordine delle sommatorie:

\begin{equation}
\begin{aligned}
\mathbb{E}[X] &= \sum_{x \in \mathcal{X}} x \sum_{m=1}^M p_{X|E_m}(x)\mathbb{P}(E_m) \\
&= \sum_{m=1}^M \mathbb{P}(E_m) \underbrace{\sum_{x \in \mathcal{X}} x p_{X|E_m}(x)}_{\mathbb{E}[X \mid E_m]}
\end{aligned}
\end{equation}

\paragraph{Esempio finale (Binomiale)}
Se applichi questa formula agli esempi precedenti ($X \sim \mathcal{B}(16, p)$), puoi ritrovare la media generale scomponendola nei due casi:

\[
\mathbb{E}[X] = \mathbb{P}(X \le 4)\mathbb{E}[X \mid X \le 4] + \mathbb{P}(X > 4)\mathbb{E}[X \mid X > 4]
\]

Oppure in egual modo:
\[
\mathbb{P}(X \in \{2,3,4,5,6\}) \mathbb{E}[X \mid X \in \{2,3,4,5,6\}] + \mathbb{P}(X \notin \{2,3,4,5,6\}) \mathbb{E}[X \mid X \notin \{2,3,4,5,6\}]
\]

Significa: La media generale dei successi (che sappiamo essere $n \cdot p = 16/3 \approx 5.33$) è uguale alla media dei successi "quando va male" (pesata per la probabilità che vada male) PIÙ la media dei successi "quando va bene" (pesata per la probabilità che vada bene).

\section{Funzioni di Variabili Aleatorie Discrete}

Sia $X$ una variabile aleatoria discreta definita su uno spazio campionario, caratterizzata da:
\begin{itemize}
    \item \textbf{Alfabeto $\mathcal{X}$:} l'insieme dei valori che $X$ può assumere.
    \item \textbf{PMF $p_X(x)$:} la funzione di massa di probabilità, dove $p_X(x) = \mathbb{P}(X = x)$.
\end{itemize}

Sia $g(\cdot)$ una funzione deterministica che opera sui punti di $\mathcal{X}$. Definiamo una nuova variabile aleatoria $Y$ come:
\[ Y = g(X) \]

L'alfabeto della nuova variabile $Y$ sarà l'insieme immagine $\mathcal{Y} = \{g(x) : x \in \mathcal{X}\}$.
\\Il problema fondamentale è ricavare la caratterizzazione statistica di $Y$ (la sua PMF e la sua media) partendo da quella di $X$. Si distinguono due casi in base alle proprietà della funzione $g$.

\subsection{Caso [a]: Funzione Biunivoca (Invertibile)}

Si verifica quando ad ogni valore $x \in \mathcal{X}$ corrisponde un unico valore $y \in \mathcal{Y}$ e viceversa ($|\mathcal{X}| = |\mathcal{Y}|$). In questo caso, esiste la funzione inversa $x = g^{-1}(y)$.

\paragraph{PMF di $Y$}
Si tratta di una semplice "ridenominazione" delle etichette dell'alfabeto. La probabilità si trasferisce direttamente dal punto $x$ al punto $y$ corrispondente:

\begin{equation}
\boxed{p_Y(y_i) = \mathbb{P}(Y = g(x_i)) = \mathbb{P}(X = g^{-1}(y_i)) = p_X(x_i)}
\end{equation}

\textit{"La probabilità di un valore nell'alfabeto di arrivo ($y_i$) è la probabilità dell'insieme di tutti i valori dell'alfabeto di partenza ($x$) che la funzione $g$ 'manda' in quel valore $y_i$."}

In questo caso però \textbf{l'insieme di tutti i valori dell'alfabeto di partenza è composto da un solo elemento}.

\subsection{Caso [b]: Funzione Non-Iniettiva}

Si verifica quando la funzione $g$ associa più valori distinti di $\mathcal{X}$ allo stesso valore di $\mathcal{Y}$ ($|\mathcal{X}| > |\mathcal{Y}|$). In questo caso, la funzione non è invertibile globalmente.

Immagina che $y_k$ sia un valore specifico nell'alfabeto di arrivo. Se la funzione non è invertibile, potrebbero esserci tanti valori di $x$ (che chiameremo $x_1^{(k)}, x_2^{(k)}, \dots, x_{L_k}^{(k)}$) che, se passati attraverso la funzione $g$, danno tutti lo stesso risultato $y_k$.

\paragraph{Esempio pratico}
Immagina $X$ (voto di un esame da 18 a 30) e la funzione $g(x)$ che trasforma il voto in "Promosso" (1) o "Bocciato" (0).
\begin{itemize}
    \item Se $Y = 1$ ("Promosso"), questo risultato deriva da molti $x$ diversi: $\{18, 19, 20, \dots, 30\}$.
    \item In questo caso, $y_k$ è "1", e i vari $x_i^{(k)}$ sono i voti dal 18 al 30.
\end{itemize}
Per un punto $y_k$ che é invertibile si applica il ragionamento di prima. 
\newpage
Per un punto $y_k$ che ha piu preimmagini avremo che:

\begin{equation}
\boxed{p_Y(y_k) = \mathbb{P}\left( \bigcup_{i=1}^{L_k} \{X = x_i^{(k)}\} \right) = \sum_{i=1}^{L_k} \mathbb{P}(X = x_i^{(k)}) = \sum_{i=1}^{L_k} p_X(x_i^{(k)})}
\end{equation}

\paragraph{notazione $x_i^{(k)}$?}
\begin{itemize}
    \item Il pedice $i$ serve a contare quanti valori di $X$ finiscono in quel punto (il primo, il secondo, ecc.).
    \item L'apice $(k)$ serve a ricordare che quei valori sono associati proprio al punto $y_k$ e non a un altro.
    \item $L_k$ è il numero totale di "punti di partenza" che finiscono in $y_k$.
\end{itemize}

Spiegazione della formula:
\begin{enumerate}
    \item $p_Y(y_k)$: È la probabilità che la nuova variabile valga $y_k$.
    \item $\mathbb{P}(\bigcup \dots)$: Dice che l'evento "$Y$ è uguale a $y_k$" accade se accade \textbf{almeno uno} degli eventi $X = x_1^{(k)}$ OPPURE $X = x_2^{(k)}$ ecc. (l'unione $\cup$ rappresenta l' "OPPURE").
    \item $\sum \mathbb{P}(\dots)$: Poiché una variabile aleatoria non può assumere due valori contemporaneamente (gli eventi sono disgiunti: o esce 18, o esce 19, non entrambi), la probabilità dell'unione diventa semplicemente la \textbf{somma} delle singole probabilità.
\end{enumerate}
Ricaptiolando:
"Nel caso [b], per calcolare la probabilità di un valore $y_k$, devo individuare tutti i valori $x$ che la funzione trasforma in quel $y_k$ e \textbf{sommare le loro probabilità originali}. Se un valore $y$ è invece 'invertibile' (proviene da un solo $x$), la sua probabilità rimane identica a quella del valore di partenza."

\subsection{Teorema Fondamentale per il Calcolo della Media}

Abbiamo visto come calcolare la PMF di una nuova variabile $Y = g(X)$. Ma se il nostro obiettivo fosse solo conoscere la sua media ($\mathbb{E}[Y]$)? Dobbiamo per forza fare tutta la fatica di calcolare prima la distribuzione $p_Y$ e poi farne la media? La risposta è no. 
\\Possiamo calcolare la media di $Y$ usando direttamente la distribuzione di $X$, senza passare per quella di $Y$. Questo risultato è noto come \textbf{Teorema Fondamentale della Media} (anche chiamato LOTUS: \textit{Law of the Unconscious Statistician}).

Seguiamo la stessa suddivisione logica vista per le PMF:

\begin{itemize}
    \item \textbf{Funzioni Biunivoche:}
    In questo caso, come abbiamo visto, c'è solo una "ridenominazione" dell'alfabeto. Ogni $y$ corrisponde a un solo $x$. Quindi sommare i valori $y$ pesati dalle loro probabilità è matematicamente identico a sommare i valori trasformati $g(x)$ pesati dalle probabilità originali:
    
    \[ \mathbb{E}[Y] = \sum_{y \in \mathcal{Y}} y p_Y(y) = \mathbb{E}[g(X)] = \sum_{x \in \mathcal{X}} g(x) p_X(x) \quad (1) \]

    \item \textbf{Funzioni Univoche (caso generale):}
    Qui più $x$ possono finire nello stesso $y$. Se usassimo la definizione classica di media su $Y$, dovremmo raggruppare questi termini.
    
    \[ \mathbb{E}[Y] = \sum_{y \in \mathcal{Y}} y p_Y(y) = \sum_{y \in \mathcal{Y}} \sum_{x: y=g(x)} g(x) p_X(x) \quad (2) \]
\end{itemize}

N.B. L'equazione (1) \textbf{include} l'equazione (2) come caso speciale.
\\Non importa se la funzione è invertibile o "molti-a-uno": la somma su tutti gli $x$ funziona sempre, perché se più $x$ vanno nello stesso $y$, la somma li includerà tutti automaticamente con i loro pesi corretti.

Adottiamo quindi la \textbf{forma generale} che prende il nome di Teorema fondamentale per il calcolo della media:

\begin{equation}
\boxed{\mathbb{E}[Y] = \mathbb{E}[g(X)] = \sum_{x \in \mathcal{X}} g(x) p_X(x)}
\end{equation}

\paragraph{Esempi}

Consideriamo una variabile aleatoria $X$ distribuita uniformemente sull'insieme (alfabeto) $\mathcal{X} = \{-2, -1, 0, 2\}$.
\\Notiamo che l'alfabeto ha cardinalità $|\mathcal{X}| = 4$.
\\Essendo uniforme, la probabilità di ogni singolo elemento è:
\[ p_X(x) = \frac{1}{4}, \quad \forall x \in \{-2, -1, 0, 2\} \]

Consideriamo la trasformazione $Y_1 = X^2$.
\\La funzione è $g(x) = x^2$.
\\\\
\textbf{Definizione del nuovo alfabeto}
L'insieme dei valori assunti da $Y_1$ è $\mathcal{Y}_1 = \{0, 1, 4\}$.
\\Osserviamo che $|\mathcal{Y}_1| = 3 < |\mathcal{X}| = 4$.
\\Siamo quindi nel \textbf{Caso [b] (Funzione non iniettiva)}: c'è stata una "collisione" di valori. I due valori originali $\{-2, 2\}$ sono collassati entrambi sul valore $4$.
\\\\
\textbf{Calcolo delle probabilità (PMF)}
Applichiamo la regola della somma per i valori collassati e la ridenominazione per gli altri.

\begin{itemize}
    \item \textbf{Per $y=0$:} Proviene solo da $x=0$.
    \[ p_{Y_1}(0) = \mathbb{P}(X=0) = \frac{1}{4} \]
    
    \item \textbf{Per $y=1$:} Proviene solo da $x=-1$.
    \[ p_{Y_1}(1) = \mathbb{P}(X=-1) = \frac{1}{4} \]
    
    \item \textbf{Per $y=4$:} Proviene sia da $x=-2$ che da $x=2$. Dobbiamo sommare le probabilità (unione di eventi disgiunti).
    \[ p_{Y_1}(4) = \mathbb{P}(\{X=-2\} \cup \{X=2\}) = \mathbb{P}(X=-2) + \mathbb{P}(X=2) = \frac{1}{4} + \frac{1}{4} = \frac{1}{2} \]
\end{itemize}

\textbf{Risultato Finale:}
\begin{equation}
\boxed{
p_{Y_1}(y) = 
\begin{cases} 
1/4 & \text{se } y = 0 \\
1/4 & \text{se } y = 1 \\
1/2 & \text{se } y = 4 \\
0 & \text{altrimenti}
\end{cases}
}
\end{equation}
\\\\
\textbf{Esempio 2:}
Consideriamo ora la trasformazione $Y_2 = 3 \sin\left(\frac{2\pi}{5}X\right)$.
\\La funzione è $g(x) = 3 \sin\left(\frac{2\pi}{5}x\right)$.

Valutiamo la funzione per ogni $x \in \mathcal{X}$:

\begin{itemize}
    \item $x = 0 \implies 3 \sin(0) = 0$
    \item $x = -1 \implies 3 \sin\left(-\frac{2\pi}{5}\right) \approx 3(-0.951) = -2.85$
    \item $x = 2 \implies 3 \sin\left(\frac{4\pi}{5}\right) \approx 3(0.588) = 1.76$
    \item $x = -2 \implies 3 \sin\left(-\frac{4\pi}{5}\right) \approx 3(-0.588) = -1.76$
\end{itemize}
\textit{(Nota: I valori numerici sono approssimati per comodità, ma il concetto importante è che sono tutti distinti).}
\\\\
\textbf{Definizione del nuovo alfabeto}
Il nuovo alfabeto è $\mathcal{Y}_2 \approx \{-2.85, -1.76, 0, 1.76\}$.
\\Osserviamo che $|\mathcal{Y}_2| = 4 = |\mathcal{X}|$.
\\Ogni valore di partenza ha generato un valore di arrivo unico e distinto. Siamo nel \textbf{Caso [a] (Funzione Biunivoca)}.
Poiché la corrispondenza è uno-a-uno, le probabilità si "trasferiscono" immutate. Se ogni $x$ aveva probabilità $1/4$, allora ogni $y$ corrispondente avrà probabilità $1/4$.

\[ p_{Y_2}(y) = \frac{1}{4} \quad \forall y \in \mathcal{Y}_2 \]

quindi: 
\begin{equation}
Y_2 \sim \mathcal{U}(\mathcal{Y}_2)
\end{equation}
